\documentclass[tikz,border=10pt]{standalone}
\usepackage{ctex}
\usepackage{amsmath}
\usetikzlibrary{calc, arrows.meta, shapes.geometric, positioning}

\begin{document}
\begin{tikzpicture}[font=\small]

% Venn 图：三个领域圆 — 向量 / 圆锥曲线 / 函数三角
% 等边三角形排列，三圆两两相交

\def\R{2.0}          % 圆半径
\def\dist{1.75}      % 圆心到中心距离

% 三圆圆心（等边三角形排列）
\coordinate (Cvec)   at (90:\dist);          % 顶部：向量
\coordinate (Cconic) at (210:\dist);         % 左下：圆锥曲线
\coordinate (Ctrig)  at (330:\dist);         % 右下：函数/三角

%──── 三个大圆（半透明填充）────────────────────
\fill[blue!18, draw=blue!60, thick]   (Cvec)   circle [radius=\R];
\fill[green!18, draw=green!60!black, thick] (Cconic) circle [radius=\R];
\fill[orange!18, draw=orange!70, thick] (Ctrig)  circle [radius=\R];

%──── 圆标签（圆心处，粗体）────────────────────
\node[blue!80!black, font=\normalsize\bfseries, align=center] at ($(Cvec)+(0,0.5)$)
  {\textbf{向量条件}\\$\vec{a},\vec{b},\vec{OA}$等};

\node[green!60!black, font=\normalsize\bfseries, align=center] at ($(Cconic)+(-0.3,-0.6)$)
  {\textbf{圆锥曲线}\\椭圆/双曲线\\抛物线/圆};

\node[orange!80!black, font=\normalsize\bfseries, align=center] at ($(Ctrig)+(0.3,-0.6)$)
  {\textbf{函数/三角}\\$\sin,\cos$\\最值/不等式};

%──── 两两交集标注────────────────────────────────
% 向量 ∩ 圆锥曲线（左上交集区域）
\node[draw=blue!40, fill=white, rounded corners=3pt,
      font=\scriptsize, align=center, inner sep=3pt]
  at ($(Cvec)!0.5!(Cconic)+(0.05,0.05)$)
  {\textbf{向量＋椭圆}\\$\vec{OA}\!\cdot\!\vec{OB}=0$\\$A,B$在椭圆上求轨迹};

% 向量 ∩ 函数/三角（右上交集区域）
\node[draw=orange!50, fill=white, rounded corners=3pt,
      font=\scriptsize, align=center, inner sep=3pt]
  at ($(Cvec)!0.5!(Ctrig)+(-0.05,0.05)$)
  {\textbf{向量＋三角}\\$\vec{a}\!\cdot\!\vec{b}=|\vec{a}||\vec{b}|\cos\theta$\\化三角最值};

% 圆锥曲线 ∩ 函数（底部交集）
\node[draw=green!50, fill=white, rounded corners=3pt,
      font=\scriptsize, align=center, inner sep=3pt]
  at ($(Cconic)!0.5!(Ctrig)+(0,-0.1)$)
  {\textbf{曲线＋代数}\\联立方程求\\最值/参数范围};

%──── 三圆公共交集（中心）────────────────────────
\node[draw=red!60, fill=red!8, rounded corners=4pt,
      font=\scriptsize\bfseries, align=center, inner sep=4pt]
  at (0,0)
  {综合压轴\\三域联动};

%──── 标题────────────────────────────────────────
\node[font=\large\bfseries] at (0,4.0) {向量综合：三域交叉 Venn 图};

\end{tikzpicture}
\end{document}
